%%%%%%%%%%%%%%%%%%%%%%%%%%%%%%%%%%%%%%%%%%%%%%%%%%%%%%%%
% 第7章 結論
%%%%%%%%%%%%%%%%%%%%%%%%%%%%%%%%%%%%%%%%%%%%%%%%%%%%%%%%
\chapter{結論}
\label{chap:conclusion}

\section{本研究のまとめ}
本研究では，動物種および撮影環境の多様性に頑健な動物行動認識の実現を目指し，
マルチモーダル表現学習と敵対的学習を組み合わせた，
ドメイン汎化に基づく新しい行動認識手法を提案した．

提案手法は，RGB 映像から抽出される外観特徴と，
オプティカルフローから抽出される運動特徴を
Gated Fusion 機構により適応的に統合する．
さらに，動物種をドメインとみなした敵対的学習を導入することで，
種固有の身体的特徴や背景バイアスに依存しない，
汎用的な行動特徴表現の獲得を実現した．

本研究によって得られた主な成果は以下の通りである：

\begin{enumerate}
  \item \textbf{マルチモーダル表現学習の有効性の実証}：
    Gated Fusion を用いることで，入力映像の特性に応じて外観と運動の重要度を動的に調整し，
    単一モダリティの限界を補完できることを実証した．
    特に，ドメインシフトの大きいターゲット環境においては，
    背景不変性の高い運動特徴が支配的な役割を果たし，認識精度の維持に大きく貢献することを明らかにした．

  \item \textbf{敵対的学習による種不変表現の獲得}：
    種分類器との敵対的学習を通じて，種情報と環境バイアスを抑制することで，
    学習データ（自然環境下のドキュメンタリー映像）とはドメインが大きく異なる未知の飼育環境に対しても有効な不変表現が獲得可能であることを示した．
    学習データにはホッキョクグマが含まれているものの，
    撮影条件や背景が劇的に変化するターゲット環境において性能を維持できた点は，
    本手法の高い環境ロバスト性を裏付けている．

  \item \textbf{実環境データでの有効性の検証}：
    動物園の固定監視カメラで撮影された，低解像度かつ低フレームレートの実映像を用いた評価実験により，
    提案手法の実運用への適用可能性を実証した．
    特に，学習データには含まれない未知の行動クラスである「常同行動」を，類似する通常の歩行動作から明確に分離・検出できたことは，
    動物福祉に貢献する実用面での大きな成果である．
\end{enumerate}

定量的評価において，提案手法はホッキョクグマ監視映像に対して
ARI 0.742，NMI 0.728 を達成し，
敵対的学習を用いない手法（ARI 0.615）と比較して
0.127 ポイントの性能向上を実現した．
この結果は，動物種をドメインとして扱う敵対的学習が
ドメイン汎化に有効であることを示している．

\vspace{1em}

結論として，本研究で提案したフレームワークは，
学習データの不足や環境の多様性という課題を抱える動物行動認識分野において，
汎用性と実用性を兼ね備えた一つの解を提示したと言える．
今後，第\ref{chap:discussion}章で議論した課題に取り組み，
本手法をさらに発展させることで，
動物園や野生環境における動物福祉の向上および保全活動に貢献する，
自律的かつ持続可能なモニタリングシステムの実現が期待される．